% SEGMENT OLP-0364-S01
% Part: lambda-calculus
% Chapter: introduction
% Section: term-revisited

\documentclass[../../../include/open-logic-section]{subfiles}

\begin{document}

\olfileid{lam}{syn}{tr}
\olsection{சொற்களை $\alpha$-சமான வகுப்புகளாகக் கருதுதல்}

இனி சொற்களை $\alpha$-சமானம்வரை கருதுவோம். அதாவது, ஒரு சொல்லை எழுதும்போது,
அது சேர்ந்துள்ள $\alpha$-சமான வகுப்பையே குறிக்கிறோம். எடுத்துக்காட்டாக,
$\lambd[a][\lambd[b][a c]]$ என்று எழுதுவது, அதற்கு $\alpha$-சமானமான
$\lambd[a][\lambd[b][a c]]$, $\lambd[b][\lambd[a][b c]]$ முதலிய எல்லாச்
சொற்களின் கணத்தையும் குறிக்கிறது.
% SOURCE CORRECTION TA-LAM-021: Tamil resolves the frozen ordinary spelling and agreement errors without changing the class construction.

மேலும், முந்தைய பகுதிகளில் $N, Q$ போன்ற எழுத்துகள் ஒரு சொல்லைக் குறிக்கப்
பயன்பட்டன; இனி அவை ஒரு வகுப்பைக் குறிக்கப் பயன்படும். தொடர்வரும் ஆய்வில்
சொற்களுக்குப் பதிலாக இந்த வகுப்புகளே நமது ஆய்வுப் பொருள்களாக இருக்கும்.
$x, y$ போன்ற எழுத்துகள் தொடர்ந்து மாறிகளைக் குறிக்கும்.

வகுப்பு~$M$ இன் ஏதேனும் ஓர் !!{element} ஐக் குறிக்க $\rep{M}$ என்ற
குறிமானத்தையும், ஒன்றுக்கு மேற்பட்டவை தேவைப்பட்டால்
$\rep{M}[0], \rep{M}[1]$ முதலியவற்றையும் பயன்படுத்துகிறோம்.

உரையை எளிமைப்படுத்த, சொற்களுக்கான குறிமானங்களை மீண்டும் பயன்படுத்துகிறோம்.
வகுப்புகளின்மீதான பின்வரும் வரையறை நம்மிடம் உள்ளது:
\begin{defn}
  \begin{enumerate}
  \item $\lambd[x][N]$ என்பது $\lambd[x][\rep{N}]$ ஐக் கொண்ட வகுப்பாக
    வரையறுக்கப்படுகிறது.
  \item $PQ$ என்பது $\rep{P}\rep{Q}$ ஐக் கொண்ட வகுப்பாக வரையறுக்கப்படுகிறது.
  \end{enumerate}
\end{defn}

$\alpha$-மாற்றம் இசைவானது என்பதால், இவை நன்கு வரையறுக்கப்பட்டவை என்பதைக்
காண்பது கடினமல்ல.

\begin{defn} \ollabel{def:fv}
  ஓர் $\alpha$-சமான வகுப்பு~$M$ இன் \emph{கட்டற்ற மாறிகள்}, அதாவது
  $\FV{M}$, என்பவை $\FV{\rep{M}}$ என வரையறுக்கப்படுகின்றன.
  % SOURCE CORRECTION TA-LAM-022: restore the FV macro omitted from the frozen formulas for class free variables.
\end{defn}

\olref[alp]{thm:fv} இல் காட்டியபடி,
$\FV{\rep{M}[0]} = \FV{\rep{M}[1]}$ என்பதால் இது நன்கு வரையறுக்கப்பட்டுள்ளது.

வகுப்புகளுக்குள் இடைநிறுத்துவதற்கான குறிமானத்தையும் மீண்டும் பயன்படுத்துகிறோம்:
\begin{defn} \ollabel{def:sub}
  $M$ இல் $y$ க்குப் பதிலாக $R$ ஐ \emph{இடைநிறுத்துவது}, அதாவது
  $\Subst{M}{R}{y}$, இடைநிறுத்தலை வரையறுக்கப்பட்டதாக ஆக்கும் ஏதேனும்
  $\rep{M}$, $\rep{R}$ க்கு $\Subst{\rep{M}}{\rep{R}}{y}$ என்று
  வரையறுக்கப்படுகிறது.
\end{defn}

\olref[alp]{cor:sub} இல் காட்டியபடி இதுவும் நன்கு வரையறுக்கப்பட்டுள்ளது.

இந்த வரையறை நமது காரணமுறையை எவ்வளவு எளிமைப்படுத்துகிறது என்பதைக் கவனிக்க.
எடுத்துக்காட்டாக:
\begin{align}
  \Subst{\lambd[x][x]}{y}{x} &\ollabel{eq:1}\\
  &= \Subst{\lambd[z][z]}{y}{x} \ollabel{eq:2}\\
  &= \lambd[z][\Subst{z}{y}{x}] \\
  &= \lambd[z][z]
\end{align}
% SOURCE CORRECTION TA-LAM-023: remove the dangling equality sign after the expression labelled eq:1; the following row supplies its equality.

இதை இன்னும் சொற்களின்மீதான இடைநிறுத்தலாகக் கருதினால் \olref{eq:1}
வரையறுக்கப்படாது; ஆனால் முன்பு கூறியபடி, இப்போது அதை வகுப்புகளின்மீதான
இடைநிறுத்தலாகக் கருதுகிறோம். அதனால்தான் \olref{eq:2} நிகழ முடிகிறது:
$\lambd[x][x]$ உம் $\lambd[z][z]$ உம் ஒரே வகுப்பைச் சேர்ந்தவை என்பதால்,
முதலாவதை இரண்டாவதால் மாற்றலாம்.

இதே காரணத்தால், இனி இடைநிறுத்தலுக்குத் தேவையான நிபந்தனைகளை நாம் தேர்ந்தெடுக்கும்
பிரதிநிதிகள் எப்போதும் நிறைவுசெய்கின்றன என்று கருதுவோம். எடுத்துக்காட்டாக,
$\Subst{\lambd[x][N]}{R}{y}$ ஐக் காணும்போது, $x \neq y$ மற்றும்
$x \notin \FV{R}$ ஆக அமையுமாறு பிரதிநிதி $\lambd[x][N]$
தேர்ந்தெடுக்கப்பட்டுள்ளதாகக் கருதுவோம்.

$\lambd[x][x]$ ஐ ஒரு ``வகுப்பு'' என்று அழைப்பது சற்றுப் புதுமையாக இருப்பதால்,
நமக்குப் பழக்கமான $\lambda$-சொற்களிலிருந்து இவற்றை வேறுபடுத்த, இனி இவற்றை
$\Lambda$-சொற்கள் (அல்லது இப்பகுதியின் மீதியில் வெறுமனே ``சொற்கள்'') என்று
அழைப்போம்.

\begin{editorial}
  சொற்களிடமிருந்து இன்னும் விடைபெற முடியாது: $\Lambda$-சொற்களின் முழு வரையறையும்
  $\lambda$-சொற்களை அடிப்படையாகக் கொண்டுள்ளது; மேலும் $\Lambda$-சொற்களின்மீது
  சார்புகளை வரையறுக்கும் ஒரு முறையை நாம் வழங்கவில்லை. அதாவது, அத்தகைய எல்லாச்
  சார்புகளும் முதலில் $\lambda$-சொற்களின்மீது வரையறுக்கப்பட்டு, இடைநிறுத்தல்களுக்குச்
  செய்ததுபோல் பின்னர் $\Lambda$-சொற்கள்மீது ``வீழ்த்தப்பட'' வேண்டும். எனினும்,
  $\Lambda$-சொற்களின்மீது சார்புகளை எவ்வாறு வரையறுக்கலாம் என்பதை வாசகர்
  உள்ளுணர்வாகப் புரிந்துகொள்ள முடியும் என்று கருதுகிறோம்.
\end{editorial}
\end{document}
